\section{Introduction}
Large language models are trained on web-scale text whose authors were rarely asked. A writer who suspects their
work was used has few options: training corpora are undisclosed, and post-hoc membership-inference attacks (MIAs)
perform close to chance on large models once members and non-members are drawn from the same distribution
\citep{duan2024membership,das2024blind}, so they cannot support an accusation
\citep{zhang2024membership}. A more promising route is for the creator to act \emph{before} publication: plant
\emph{canaries}, unique strings with no reason to appear elsewhere, and later test whether a model has memorized
them \citep{carlini2019secret,meeus2024copyright,wei2024proving}.

The practical question is quantitative. A writer controls how many canaries they plant and what they look like,
but not how the developer's pipeline treats them: deduplication can remove repeats, Unicode normalization can
erase invisible marks, and paraphrasing can destroy verbatim strings. Nor does the writer control the access they
get: commercial models usually expose text, not log-probabilities. We therefore ask:
\textbf{how many copies of which canary does a writer need before training on their work becomes provable, and
with what access?}

We answer with a controlled study in which ground truth is known. We continue pretraining \pythia{160M} and \pythia{410M}
\citep{biderman2023pythia} on Wikipedia into which we plant canaries of four types at four copy counts; matched
canaries from the same generators that are never planted serve as non-members, and the member/non-member split
is randomized after generation, which makes a rank test of member vs.\ non-member scores a valid permutation-style
test with exact false-positive control. Our contributions:
\begin{itemize}
  \item A \textbf{canary scorecard} (Table~\ref{tab:scorecard}): for each canary type, the minimum number of copies
        detectable at $p<0.01$, whether it survives deduplication, normalization and paraphrasing, whether
        text-only access suffices, and its false-positive rate on a model that never saw it.
  \item \textbf{Dose--response and survival measurements} under realistic pipelines and two access levels
        (Figures~\ref{fig:dose} and~\ref{fig:survival}) at two model sizes, including the finding that single copies are only detectable when repeated canaries of the same kind are also present.
  \item A \textbf{pre-registered audit protocol} whose false-positive control we verify empirically on the clean
        base models (Appendix~\ref{app:qq}), and released code that turns the protocol into a tool for creators.
\end{itemize}
